\section{When is $(z,K)$ a sufficient state?}\label{sec:aggregate}

Households need only prices, and prices depend on the path of capital and labour. So the right
question is not whether $(z,K)$ determines the distribution, but whether it determines that path.

\begin{definition}
A sequence-state equilibrium $\mu^*$ satisfies \emph{$(z,K)$-sufficiency} if any two histories with
the same current shock and the same capital have the same capital and labour at every node of their
future trees.
\end{definition}

\noindent
A \emph{recursive equilibrium on $(z,K)$} consists of capital sets $\mathcal K_z$, a law of motion
$K'=\Psi(z,K,z')$ keeping capital in the sets, labour $L(z,K)$, and value and saving functions of
$(z,K,a,e)$ satisfying the Bellman equation with continuation at $(z',\Psi(z,K,z'))$, together with
consistency of aggregation: each $(z,K)$ carries a nonempty set of distributions with mean $K$ and
labour $L(z,K)$, mapped by the saving policy and the employment transitions into the set at
$(z',\Psi(z,K,z'))$. It \emph{represents} $\mu^*$ if at every history the capital lies in
$\mathcal K_z$, the distribution lies in the set at $(z,K)$, and households save by the $(z,K)$
policy.

\begin{theorem}\label{thm:iff}
For every sequence-state equilibrium, $(z,K)$-sufficiency holds if and only if some recursive
equilibrium on $(z,K)$ represents it. For small capital shares (as in Theorem~\ref{thm:stationary}):
when $(z,K)$-sufficiency holds, the capital sets are compact, the law of motion $\Psi$ is continuous
on them, and the law of $(z,K)$ is the unique invariant distribution of the $(z,K)$ process on the
capital sets, and it is ergodic.
\end{theorem}
\formal{kSufficient\_iff\_exists\_represents, minMPC\_smallShare\_aggregateState, stationary\_aggLaw}

\noindent
Continuity needs no injectivity: capital is a quotient map from the compact set of histories with a
given shock. Uniqueness of the stationary law needs no Lipschitz test functions and no measurable
selection of distributions: bounded continuous test functions are uniformly continuous on the
compact capital sets, and capital forgets old shocks.

\begin{theorem}\label{thm:repagent}
With a single idiosyncratic state, every household holds aggregate capital in the small-share
equilibrium, so capital determines the distribution and $(z,K)$-sufficiency holds: the stochastic
growth model of \citet{brockmirman1972} is recovered as a recursive equilibrium on $(z,K)$ with a
unique stationary distribution.
\end{theorem}
\formal{minMPC\_smallShare\_repAgent}

\paragraph{An order route.} If the cross-sections the equilibrium visits at each shock form a chain
for first-order stochastic dominance in assets, then capital determines the cross-section, because
two comparable cross-sections with the same mean are equal. The chain property follows from two
conditions: the law of motion preserves dominance, and after ordered shocks the next cross-sections
are ordered whatever the pasts. The first holds, by Tarski comparative statics, when saving rises
with the capital path (Theorem~\ref{thm:monotone}); the second reduces to finitely many comparisons
between extreme histories.

\begin{theorem}\label{thm:order}
Dominance with equal means is equality. If saving rises with the capital path, sequential
equilibrium is unique, the equilibrium is continuous in the history, a better newest shock raises
the cross-section, and the extreme comparisons hold, then $(z,K)$-sufficiency holds.
\end{theorem}
\formal{eq\_of\_dominates\_of\_mean\_eq, kInjective\_of\_chain, kSufficient\_of\_savingMonotone,
kSufficient\_of\_extremes}

\noindent
The extreme comparisons say that one period of a better shock outweighs a permanent boom against a
permanent bust. They force capital to encode the whole shock history, which needs the economy to
forget at least half of a capital deviation each period with two shocks. They cannot hold at the
persistence of Krusell and Smith's calibration, which says nothing about $(z,K)$-sufficiency there,
only that this sufficient condition is unavailable.
